\documentclass[11pt]{article}
\usepackage[a4paper,margin=25mm]{geometry}
\usepackage{amsmath,amssymb,amsthm,booktabs,hyperref}
\newtheorem{theorem}{Theorem}
\newtheorem{lemma}{Lemma}
\newcommand{\D}{\mathcal D}
\newcommand{\C}{\mathcal C}
\newcommand{\E}{\mathcal E}
\newcommand{\FP}{\operatorname{FPdim}}
\title{A six-versus-four Galois obstruction for $R6_{9,3748}$}
\author{}\date{1 October 2026}
\begin{document}\maketitle

\begin{theorem}
The based ring $R6_{9,3748}$ admits no fusion categorification over an
algebraically closed field of characteristic zero. In particular it has no
pivotal, spherical, pseudo-unitary, unitary, or braided categorification.
\end{theorem}

The accompanying \texttt{structural\_reduction.tex} proves two categorical
reductions without a pivotality assumption: every hypothetical categorification
is spherical and Galois conjugate to a pseudo-unitary one; and it admits a
fully faithful central lift of its $\operatorname{PSU}(2)_5$ subcategory,
with a nondegenerate induced braiding. Its nine-case sphericalization
certificate and the original fifteen-case relative-center certificate are
part of the proof. This note excludes the resulting pseudo-unitary category.
Only the braided Drinfeld center is factored below; the ambient fusion
category is not assumed braided.

\paragraph{Dimension field and center factorization.}
Write the ring as
\[
 A[X]/(X^2-1-QX),\qquad
 a^2=1+b,\quad ab=a+b,\quad b^2=1+a+b,\quad Q=2+2a+b,
\]
with ordered basis $(1,a,b,X,aX,bX)$ corresponding to source labels
$(0,1,5,2,3,4)$. Henceforth the same lower-case letters denote positive
dimensions:
\[
 a=2\cos(\pi/7),\quad b=a^2-1,\quad h=\sqrt{13},\quad
 q=2+2a+b,\quad x=(q+bh)/2.
\]
The identities
\[
 a^3-a^2-2a+1=0,\qquad q^2+4=13b^2,\qquad x^2-qx-1=0
\]
hold exactly. Set
\[
 d=1+a^2+b^2,\qquad f=d(1+x^2)=\FP(\C).
\]
The six numbers $1,a,b,x,ax,bx$ are linearly independent over $\mathbb Q$;
thus the dimension homomorphism $\mathbb Q\otimes K_0(\C)\to
L=\mathbb Q(a,h)$ is injective. In particular a dimension determines
the underlying Grothendieck class of an object of $\C$.

By the central lift and the nondegenerate-subcategory decomposition theorem
\cite[Theorem 3.13]{DGNO},
\[
 Z(\C)\simeq\D\boxtimes\E,\qquad
 \dim\E=f^2/d,
\]
where $\E$ is the modular centralizer. All dimensions are positive: the
canonical spherical structure of the pseudo-unitary center restricts to its
two factors. The forgetful functor $F:Z(\C)\to\C$ restricts on $\D$
to the given inclusion. We use the same letter $F$ on $\E$.

\paragraph{The six known simples of $\E$.}
Let $I$ be right adjoint to $F$, and let $n$ be the multiplicity vector of
$I(1)$. Since $K_0(\C)$ is commutative of rank six, the formal-codegree
description of induction gives six distinct simple summands, each of
multiplicity one \cite[Theorem 2.13]{Ostrik}. Their twists are all one
\cite[Theorem 2.5]{Ostrik}. For normalized modular $S$ in the center,
\[
 Sn=n
\]
by \cite[Proposition 2.8]{Ostrik}. Therefore, for every Galois
automorphism $\sigma$, the signed permutation
$G_\sigma=\sigma(S)S^{-1}$ fixes $n$. Its underlying permutation preserves
the support of $I(1)$, which contains the unit.

For clarity, the following standard modular Galois identities are used:
\begin{align}
 d_{\widehat\sigma(i)}
   &=d_{\widehat\sigma(0)}\,|\sigma(d_i)|,
   \label{eq:dim-galois}\\
 \theta_{\widehat\sigma(i)}
   &=\theta_{\widehat\sigma(0)}\,\sigma^2(\theta_i).
   \label{eq:twist-galois}
\end{align}
The first follows from the signed action on normalized $S$, using positive
dimensions. For the second take a modular representation with normalized
$t_i=\theta_i/y$; the theorem $t_{\widehat\sigma(i)}=\sigma^2(t_i)$
of \cite[Theorem II(iii)]{DLN}, divided by its case $i=0$, gives precisely
\eqref{eq:twist-galois}. This unit-twist factor must not be omitted in
general.

In $Z(\C)$, the unit's Galois images have dimensions $f/\sigma(f)$.
The formal-codegree polynomial
\[
 z^6-637z^5+22295z^4-298116z^3+1854944z^2-5274997z+5274997
\]
is irreducible and has six distinct positive roots. Thus these six images
are distinct and exhaust the summands of $I(1)$.

The normalized $S$-matrix of the Deligne product is the tensor product
of the two normalized $S$-matrices. Consequently its Galois permutation
sends $(1,1)$ to
$(\widehat\sigma_{\D}(1),\widehat\sigma_{\E}(1))$.
This directly identifies the factors of the six induced simples; no
assumption about restriction of fusion characters is required.
The three embeddings of $\mathbb Q(a)$ send $a$ to
$a,a^2-a-1,2-a^2$. Since
$d_{\widehat\sigma_{\D}(1)}^2=d/\sigma(d)$, the respective simple
labels in $\D$ are $1,b,a$.
The $\E$ factor therefore has dimension
$(f/\sigma(f))/d_{\widehat\sigma_{\D}(1)}$.

Using the injective dimension map gives the following exact forgetful
vectors $v_e=[F(e)]$, in the six-element basis above. The verifier checks
every identity in $L$.
\begin{center}
\begin{tabular}{cccl}\toprule
$\sigma(a)$&$\sigma(h)$&$\D$ label&$v_e$\\\midrule
$a$&$h$&$1$&$(1,0,0,0,0,0)$\\
$a$&$-h$&$1$&$(1,0,0,2,2,1)$\\
$a^2-a-1$&$h$&$b$&$(0,0,1,1,3,4)$\\
$a^2-a-1$&$-h$&$b$&$(0,0,1,0,0,1)$\\
$2-a^2$&$h$&$a$&$(0,1,0,1,2,1)$\\
$2-a^2$&$-h$&$a$&$(0,1,0,1,1,2)$\\\bottomrule
\end{tabular}
\end{center}
These are six distinct $\E$ simples, with distinct dimensions. Their twists
are inverse to the twists of their $\D$ labels because the corresponding
product objects lie in $I(1)$. Since $N_{\mathbb Q(a)/\mathbb Q}(d)=49$, the Cauchy theorem
\cite[Theorem 3.9]{BNRW} applied to $\D$ shows that all six twists
have seven-primary order.
In particular the second object, denoted $e_-$, has
\begin{equation}
 d_{e_-}=x^2,\qquad \theta_{e_-}=1.
 \label{eq:eminus}
\end{equation}

\paragraph{The induction Gram budget.}
Let $N_i$ be the six fusion matrices of $\C$. The induction formula
$FI(Y)=\bigoplus_j jYj^*$ \cite[Proposition 2.4]{Ostrik} and adjunction
give
\[
 \sum_{z\in\operatorname{Irr}Z(\C)}[F(z)][F(z)]^t
 =G:=\sum_{i=0}^5N_i^2
 =\begin{pmatrix}
6&2&4&10&15&16\\2&10&6&15&26&31\\4&6&12&16&31&41\\
10&15&16&72&115&139\\15&26&31&115&211&254\\16&31&41&139&254&326
\end{pmatrix}.
\]
Every simple $e\in\E$ supplies the three simple objects
$1\boxtimes e,a\boxtimes e,b\boxtimes e$ with forgetful vectors
$v_e,N_av_e,N_bv_e$. Subtracting the eighteen vectors from the six known
objects above leaves
\[
 G-\sum_{e\ \mathrm{known}}\sum_{i=0}^2(N_iv_e)(N_iv_e)^t
 =\begin{pmatrix}0&0\\0&T\end{pmatrix},\qquad
 T=\begin{pmatrix}24&36&44\\36&68&80\\44&80&104\end{pmatrix}.
\]
All contributions are nonnegative outer products. Hence every further
simple $y\in\E$ has
\[
 F(y)=XU,\qquad U=u_0+u_1a+u_2b,
 \quad u=(u_0,u_1,u_2)^t\in\mathbb Z_{\geq0}^3\setminus\{0\}.
\]
None has the same dimension as one of the six known objects, since their
forgetful vectors have a nonzero component in $A$.
Let $M_0,M_a,M_b$ be the rank-three fusion matrices of $A$. The entire
$\D$ block of $y$ contributes
$\sum_i(M_iu)(M_iu)^t$ to $T$. Its upper-left entry is
$u_0^2+u_1^2+u_2^2$. Therefore
\begin{equation}
 \sum_{y\ \mathrm{unknown}}\|u_y\|^2=24.
 \label{eq:budget}
\end{equation}

\paragraph{Cubic Galois action bounds dimension multiplicities by four.}
Take a cyclotomic Galois automorphism $\tau$ restricting to
$a\mapsto a^2-a-1$ and fixing $h$. For any unknown simple, applying
\eqref{eq:dim-galois} and the known Galois image of the unit gives
\[
 d_{\widehat\tau(y)}=x\left|
 \frac{b^2}{\tau(b)}\tau(U)\right|.
\]
In the basis $(1,a,b)$, multiplication of $\tau(U)$ by
$b^2/\tau(b)$ is represented by the exact matrix
\[
 H=\begin{pmatrix}0&-1&1\\-1&-1&1\\-2&0&1\end{pmatrix},
 \qquad H^3=-I.
\]
The forgetful vector of $\widehat\tau(y)$ must be either $Hu$ or $-Hu$,
whichever has nonnegative entries. A vector with mixed signs would be
impossible by injectivity of the dimension map. In particular unknown
simples remain unknown. Their dimensions form orbits of size one or three;
the number of simple objects of each dimension is constant along each
orbit, because the Galois permutation is bijective.

Every such dimension orbit has total squared coordinate norm at least six.
A fixed positive vector satisfies $Hu=-u$, and the positive integral
solutions are $u=n(1,2,1)$, $n\geq1$, of norm squared $6n^2$.
For a three-element orbit, a total smaller than six would require a vector
of norm squared one. The only such nonnegative integral vectors are the
three coordinate unit vectors. Their complete positive-sign $H$ orbits are
\[
\begin{array}{c|c}
 \text{orbit}&\text{sum of squared norms}\\\hline
 (1,0,0),(0,1,2),(1,1,2)&12\\
 (0,1,0),(1,1,0),(1,2,2)&12\\
 (0,0,1),(1,1,1),(0,1,1)&6.
\end{array}
\]
Thus if $m$ distinct simples share any unknown dimension, their complete
dimension orbit contributes at least $6m$ to \eqref{eq:budget}; consequently
\begin{equation}
 m\leq4.
 \label{eq:four}
\end{equation}
The known dimensions have multiplicity one, so the same bound holds for
every simple dimension in $\E$. This argument uses no enumeration of
possible modular rings, and no simple-dimension $d$-number filter.

\paragraph{A required twist forces multiplicity at least six.}
The exact norm is
\[
 N_{L/\mathbb Q}(\dim\E)=N_{L/\mathbb Q}(f^2/d)=7^4\,13^6.
\]
The Cauchy theorem for spherical fusion categories
\cite[Theorem 3.9]{BNRW} implies that thirteen divides the
Frobenius--Schur exponent of $\E$. For a modular category this exponent
is the order of its twist matrix. Hence some simple $y$ has twist with a
nontrivial thirteen-primary component. Such a $y$ is unknown, because the
six known twists have seven-primary order.

Choose a Galois automorphism $\sigma$ which fixes all prime-to-thirteen
roots of unity in a common cyclotomic field for the modular representations
and acts by exponent two on its thirteen-primary roots. The Chinese
remainder theorem permits this choice. Then
\[
 \sigma(a)=a,\qquad \sigma(h)=-h,\qquad\sigma(x)=-1/x,
\]
since two is a quadratic nonresidue modulo thirteen. Its Galois image of
the unit in $\E$ is the known object $e_-$ of \eqref{eq:eminus}.
For every unknown simple of dimension $xU$, equations
\eqref{eq:dim-galois}--\eqref{eq:eminus} now give
\[
 d_{\widehat\sigma(y)}=x^2\,|\sigma(xU)|=xU,
 \qquad
 \theta_{\widehat\sigma(y)}=\sigma^2(\theta_y).
\]
Thus $\widehat\sigma$ preserves each equal-dimension block and acts on
the thirteen-primary part of its twists by exponent four. The order of
four modulo thirteen is six:
\[
 1,4,3,12,9,10,1.
\]
For any nontrivial root of order $13^j$, the corresponding exponent-four
orbit has length at least six, already by reduction modulo thirteen.
It follows that the block containing $y$ has at least six distinct simples.
This contradicts \eqref{eq:four}, excluding the pseudo-unitary category.
The spherical/Galois reduction excludes every characteristic-zero
categorification, proving the theorem.

\paragraph{Exact verification and dependencies.}
The short verifier \texttt{verify\_center\_exclusion.py} checks the source
tensor, the dimension-field identities and injectivity, the six rows,
their Galois $\D$ labels, the full induction Gram subtraction, the cubic
matrix and its dimension multiplier, the fixed ray and all three
unit-vector orbits, the global norm, quadratic conjugation, and the
order-six residue orbit. Run with \texttt{--verify} to replay
\texttt{center\_exclusion\_certificate.json}; it also succeeds under
\texttt{python -O}. It does not claim that arithmetic alone proves the
categorical reductions.
The independent \texttt{verify\_sphericality.py --verify} replays the
nine sphericalizations. The standalone \texttt{verify\_relative\_center.py --verify} replays all
fifteen rank-eighteen candidates. All three verifiers use the accompanying
complete frozen source row, \texttt{input\_tensor.json}, whose SHA256 is
pinned in the helper verifier. No former research directory or external
source checkout is required. Source provenance is recorded in
\texttt{source\_metadata.json}; \texttt{verify\_all.py} runs all three
verifiers under \texttt{python -O}.

\begin{thebibliography}{9}
\bibitem{Ostrik} V.~Ostrik, \emph{Pivotal fusion categories of rank 3},
\href{https://arxiv.org/abs/1309.4822}{arXiv:1309.4822},
Propositions 2.4 and 2.8, Theorems 2.5 and 2.13.
\bibitem{DGNO} V.~Drinfeld, S.~Gelaki, D.~Nikshych, V.~Ostrik,
\emph{On braided fusion categories I},
\href{https://arxiv.org/abs/0906.0620}{arXiv:0906.0620}, Theorem 3.13.
\bibitem{DLN} C.~Dong, X.~Lin, S.-H.~Ng,
\emph{Congruence property in conformal field theory},
\href{https://arxiv.org/abs/1201.6644}{arXiv:1201.6644}, Theorem II.
\bibitem{BNRW} P.~Bruillard, S.-H.~Ng, E.~C.~Rowell, Z.~Wang,
\emph{Rank-finiteness for modular categories},
\href{https://arxiv.org/abs/1310.7050}{arXiv:1310.7050}, Theorem 3.9.
\end{thebibliography}
\end{document}
